% Benchmark chapter: experimental design (paper section) and the locked-set results in thesis length. Result tables are the generated files in
% ../paper/tables; numbers in the text are copied from results/final/stats_held.json. The paper's results.tex is NOT input (labels would clash).
\section{Experiments}
\label{sec:experiments}
% PLAN (about 2 pages and one table). Design only: no result appears here. Everything below restates SAP.md / SAP_v2.md, which were frozen before the
% locked evaluations (tools/freeze.py, tools/freeze_v2.py); if this text and SAP_v2.md differ, SAP_v2.md wins and this text is corrected.
%   6.1 Models and comparators (ours, no-twin control, v1, Qure released weights, identity, CLAHE, homomorphic filtering).
%   6.2 Arms (Table arms) and the dose-response over the twin weight.
%   6.3 Protocols per hypothesis H1'-H5' (what is computed, on which set, with which test).
%   6.4 Metrics. 6.5 Downstream protocol. 6.6 Encoder screen (selection rule fixed before the screen).
\subsection{Models and comparators}
\emph{Ours} is the ensemble of three ConvNeXt-V2-T UNet++/CBAM suppressors (seeds 0--2) trained on the Qure projections with twin weight $w=16$
(Section~\ref{sec:method_v2}). It replaced the SwinV2-T configuration of the analysis plan under the pre-declared replacement rule (R2) before the
freeze: ConvNeXt-V2-T and RAD-DINO both met it on development data, and the pre-declared tie-break (mean absorption within 0.005, then the
higher bone-region PSNR) chose ConvNeXt-V2-T at two seeds each. With the third seeds the absorption gap is 0.0055, just outside the band; the
choice was not reopened, and both encoders are reported. The SwinV2-T and RAD-DINO pairs (twin and no twin, three seeds each) are secondary.
Comparators are: the \emph{no-twin control} (same encoder, data and recipe, $w=0$, three seeds); the v1 model trained on LIDC-IDRI only (\texttt{clia},
pre-specified pick of the v1 study); the publicly released weights of \citet{angaitkar2026anatomy} (SwinV2-B at $1024$~px; research licence);
and three controls that need no training: identity, CLAHE \citep{zuiderveld1994clahe} and a homomorphic filter
\citep{oppenheim1968homomorphic}. All methods are scored with the same code on the same images.

\subsection{Arms}
Table~\ref{tab:arms} lists the arms. Each arm changes one factor of its parent. The confirmatory dose-response varies the twin weight over
$w\in\{0,4,8,16,32,64\}$ in the SwinV2-T family with 3, 2, 2, 4, 2 and 2 seeds ($w=1$ appears only in the screening). The insertion-only arm and the
last four rows form the ablation of the regulariser (two seeds each, SwinV2-T, $w=16$).

\begin{table}[htbp]
\centering\small
\renewcommand{\arraystretch}{1.15}
\caption{Training arms. ``Lesions'' is the maximum number of lesions inserted into the base image before the physics (counterfactual insertion);
``Twin'' is the twin weight $w$, varied over $\{4,8,16,32,64\}$ in the \texttt{qp\_twin}$w$ arms ($w=1$ in the screening only). The parent of an arm is \texttt{recipe} for the
LIDC arms, \texttt{qp} for the first Qure arms and \texttt{qp\_twin16} for the last four rows.}
\label{tab:arms}
\begin{tabular}{@{}p{0.25\linewidth}p{0.08\linewidth}p{0.09\linewidth}p{0.08\linewidth}p{0.26\linewidth}@{}}
\toprule
Arm & Source & Lesions & Twin & Change \\
\midrule
\texttt{recipe} & LIDC & 0 & none & v1 baseline \\
\texttt{nodule\_loss} & LIDC & 0 & none & lesion-weighted loss \\
\texttt{clia} & LIDC & 3 & none & counterfactual insertion \\
\texttt{clia\_nodule} & LIDC & 3 & none & insertion and lesion-weighted loss \\
\texttt{qp} & Qure & 0 & none & Qure projections (no-twin control) \\
\texttt{qp\_clia} & Qure & 3 & none & insertion only, no tie term \\
\texttt{qp\_twin}$w$ & Qure & 3 & varied & twin regulariser \\
\texttt{qp\_twin16\_sg} & Qure & 3 & 16 & stop-gradient on the un-pasted branch \\
\texttt{qp\_twin16\_wide} & Qure & 3 & 16 & support dilated with a $31\times31$ window \\
\texttt{qp\_twin16\_all} & Qure & 3 & 16 & whole-image support \\
\texttt{qp\_twin16\_n4} & Qure & 3 & 16 & four lesions per twin \\
\bottomrule
\end{tabular}
\end{table}

\subsection{Protocols per hypothesis}
The analysis plan (SAP.md for the v1 family, SAP\_v2.md for v2) was frozen before any held-out or locked evaluation, on 1~November 2026
(\path{freeze_v2.json}: hashes of both plans, the SZCH development and held-out lists, the split files and the 52 weight files, at commit
\texttt{dd52b6a}). Paired differences are tested with the Wilcoxon signed-rank test \citep{wilcoxon1945individual} and given percentile bootstrap intervals
\citep{efron1993bootstrap}. Multiplicity is handled by the Holm procedure \citep{holm1979simple} on the $p$-values within each family; confidence intervals are unadjusted
95\,\% intervals; effect sizes are always reported with them, and a hypothesis that fails is reported as failed. The external validation on
BDCXR-3257 has its own addendum, frozen before any model saw those images (Section~\ref{sec:bdcxr_protocol}).
\begin{description}
\item[H1$'$, dose-response.] On VinDr-RibCXR (245 images, six sites each, lesions from LIDC test CTs, $\alpha=1$), the per-site absorption $A$ averaged over
 seeds falls with the twin weight: the slope of $A$ on $\log_2(1+w)$, with an image-cluster bootstrap interval, passes if its upper bound is below zero.
 The same analysis on the Qure-projection test split (exact truth) is secondary.
\item[H2$'$, fidelity.] On the 51 held-out SZCH pairs, bone-region PSNR of ours against the reference model is non-inferior within a margin of 0.5~dB
 (one-sided 95\,\% lower bound of the paired mean difference, image-level bootstrap); SSIM has a margin of 0.01 and is reported either way; ours against the
 no-twin control uses a margin of 0.25~dB (the cost of the regulariser). Metrics are computed after a per-image gain and offset fit to the DES reference
 (``calibrated'', primary) and without it (``raw'').
\item[H3$'$, lesion absorption.] On VinDr-RibCXR sites, $A$(ours)$<A$(reference model) and $A$(ours)$<A$(no-twin control), as paired site-level
 differences with image-cluster bootstrap intervals, Holm over the two contrasts, and an interaction between sites that overlap an annotated rib and sites
 that do not.
\item[H4$'$, evidence without insertion.] (a) Contrast-to-noise change of real nodules on JSRT and Node21, ours against the reference model and against
 the no-twin control, for all lesions and for lesions with at least 15\,\% bone overlap. (b) Downstream: NIH ChestX-ray14 macro AUROC per input
 (original, ours, reference model, CLAHE, original plus ours; three seeds); TBX11K, Node21 and VinBigData detection by FROC and mAP@50; recall at two false
 positives per image for lesions with and without bone overlap, as a difference in differences against the original input.
\item[H5$'$, reliability.] Spearman correlation between ensemble spread and absolute bone error on the SZCH held-out pairs (pseudo-truth bone from the
 least-squares-aligned DES image); split-conformal coverage calibrated on Qure-projection validation CTs and measured on the Qure-projection test split
 (nominal 90\,\%) and, as measured under shift, on SZCH.
\end{description}

\subsection{Metrics and downstream models}
Paired fidelity uses PSNR and SSIM \citep{wang2004ssim} on the bone region and on the whole image, band-limited SSIM, MS-SSIM \citep{wang2003msssim}, LPIPS \citep{zhang2018lpips}
and a high-frequency energy ratio; the same smoothing is applied to every method and never tuned. The rib benchmark reports the relative reduction of
Sobel gradient energy on the rib-boundary band inside the lungs and on lung pixels away from ribs, their difference (selectivity) and the ROC-AUC of the
gradient-reduction map. Detection is scored by FROC sensitivity at fixed false-positive rates, averaged as the competition performance metric
\citep{vanginneken2010anode09,niemeijer2011cpm}, and by mAP@50; AUCs are compared with the paired DeLong test \citep{delong1988auc}. The downstream classifier is a DenseNet-121
\citep{huang2017densenet} and the detector a Faster R-CNN with a ResNet-50 FPN backbone \citep{ren2017fasterrcnn}, with identical settings for every
input variant (\path{s4/cls.py}, \path{s4/det.py}). The classifier sees $256$~px images, is trained for 10 epochs with AdamW
\citep{loshchilov2019adamw} (learning rate $3\times10^{-4}$, weight decay $10^{-4}$, batch 64, cosine schedule) and keeps the epoch with the
best validation macro-AUROC. The detector sees $512$~px images, is trained for 12 epochs with SGD (learning rate 0.01, momentum 0.9, weight
decay $10^{-4}$, batch 8) and keeps the epoch with the best validation CPM. A multi-channel input (original plus soft, or original, soft and bone)
stacks the images as channels. Detection metrics treat tied scores as one threshold, so they do not depend on the order in which a sort
returns ties.

\subsection{Encoder screen}
To check whether a newer encoder helps under the same recipe, ConvNeXt V2 \citep{woo2023convnextv2}, DINOv3 \citep{simeoni2025dinov3} and RAD-DINO
\citep{perezgarcia2025raddino} were trained at a reduced budget under the same recipe, and an encoder moves to the full budget only if
it meets a rule fixed in advance (RUNPLAN.md, rule R1): at \num{100000} samples, absorption $A\le0.070$ on the development audit with
bone-region PSNR$_{cal}\ge21.60$~dB on the development pairs, or PSNR$_{cal}\ge21.95$~dB with $A\le0.100$. The thresholds come from the
SwinV2-T range at that budget and were written after one SwinV2-B run had been seen; the decision at the full budget used rule R2 (see Models and comparators above). Pretrained encoders other than SwinV2 are trained with a learning rate $0.25\times$ that of the decoder;
ConvNeXt-V2-T diverged at the full rate (fp16 overflow) and was rerun at $0.25\times$, an extension of the rule written for vision transformers.
DINOv3 ViT-S/16 did not improve on SwinV2-T, RAD-DINO ViT-B/14 missed R1 on PSNR by 0.05~dB but was continued at the full budget as an
exploratory encoder, and ConvNeXt-V2-T met R1.

\subsection{External validation protocol (BDCXR-3257)}
\label{sec:bdcxr_protocol}
An addendum to the analysis plan (SAP\_bdcxr.md) was frozen before any model of this study saw a BDCXR-3257 image; its one change, the
duplicate grouping (Section~\ref{sec:data}), was made from the images alone before any model ran and is recorded in the deviations log. All
suppressors are the frozen models; only the downstream classifiers are trained. (E1, primary) Pneumonia classification with three one-vs-rest
outputs, inputs as in H4$'$b (original, ours, reference model, CLAHE, original plus ours), the classifier settings of the downstream protocol,
five folds stratified by class over the near-duplicate groups (fold seed 0) with 10\,\% of each training part, by whole groups, for epoch selection,
and three seeds per fold and input (75 runs). Out-of-fold probabilities, averaged over seeds, give one prediction per image. The primary measure
is the pooled macro-AUROC; ours against the original image is non-inferior if the lower bound of the 95\,\% interval of the difference (group
bootstrap) exceeds $-0.01$, and superior if it exceeds 0; per class, paired DeLong tests with Holm over the three classes. (E2, secondary) The
lesion audit of Section~\ref{sec:audit} on a fixed class-stratified subset of 300 images (seed 7), six sites per image, with bone overlap from the
reference model's bone map; images whose lung mask covers less than 2\,\% of the image are skipped and counted (none was). (E3, descriptive) The
per-image ensemble spread over the lung field, compared with the held-out SZCH pairs and correlated with age.


\section{Fidelity on held-out paired radiographs}
On the 51 held-out SZCH pairs (Tables~\ref{tab:szch_heldout} and~\ref{tab:h2}) the bone-region PSNR$_{cal}$ of ours is 21.46\,dB and that of
the reference model 21.51\,dB; the paired difference is $-0.06$\,dB with a one-sided 95\,\% lower bound of $-0.15$\,dB, well inside the
non-inferiority margin of $-0.5$\,dB. Against the no-twin control the difference is $+0.01$\,dB (lower bound $-0.02$\,dB against a margin of
$-0.25$\,dB). Both parts of H2$'$ that concern bone-region fidelity pass. SSIM$_{cal}$ fails: ours is lower than the reference model by 0.021,
with a lower bound of $-0.023$ against a margin of $-0.01$, as the plan had anticipated from the development data.

The SSIM difference was examined on development data before the freeze (Appendix~\ref{app:negative}). The dual-energy soft-tissue image has
less high-frequency energy than a radiograph, so full-resolution SSIM against it rewards smooth outputs. On the held-out pairs, ours keeps
about 9\,\% more high-frequency energy than the dual-energy image in non-bone regions, the reference model about 1\,\% less; with a Gaussian
band limit of $\sigma=4$ the SSIM of the two models is 0.946 and 0.948, so almost all of the gap is fine texture. Whether a $1024$~px output,
as the reference model has, would close the rest was not tested. LPIPS points the same way (0.234 for ours, 0.166 for the reference model).
BS-LDM, trained on dual-energy pairs from the same hospital as the test images, reaches 28.57\,dB bone-region PSNR$_{cal}$; it is shown
for orientation only.
\begin{table}[htbp]
\centering\small
\caption{Held-out SZCH-X-Rays pairs: bone-region and whole-image metrics for ours and the comparators, calibrated (primary) and raw.}
\label{tab:szch_heldout}
\fittable{tables/szch_heldout.tex}
\end{table}
\begin{table}[htbp]
\centering\small
\caption{H2$'$ non-inferiority on the 51 held-out SZCH-X-Rays pairs (paired differences, image bootstrap, one-sided 95\,\% lower bound).}
\label{tab:h2}
\fittable{tables/h2_noninferiority.tex}
\end{table}
\begin{figure}[htbp]
\centering
\safefig{figures/pareto.pdf}
\caption{Lesion absorption on VinDr-RibCXR against bone-region PSNR$_{cal}$ on the 51 held-out SZCH-X-Rays pairs: the SwinV2-T family at
twin weights $0\ldots64$, ours and the no-twin control (ConvNeXt-V2-T, $K=3$), the reference model, the v1 model and the identity.}
\label{fig:pareto}
\end{figure}
\begin{figure}[htbp]
\centering
\safefig{figures/qualitative_szch.png}
\caption{Held-out SZCH-X-Rays cases (CC BY-NC 4.0, \citealp{sun2026bsldm}), chosen by rule: the held-out images at the 25th, 50th and 75th
percentile of the bone-region PSNR$_{cal}$ of ours. Radiograph, dual-energy soft-tissue reference, the soft-tissue estimates of ours and of the
reference model, and the bone estimate of ours.}
\label{fig:qualitative}
\end{figure}

\section{Rib benchmark}
On VinDr-RibCXR the rib benchmark measures how much gradient energy a suppressor removes on rib-boundary pixels and on other lung pixels
(Table~\ref{tab:rib_benchmark}). The reference model removes the most on both (48\,\% and 22\,\%), ours less (30\,\% and 8\,\%), and the
no-twin control lies in between (37\,\% and 14\,\%). Selectivity, the difference of the two, is 0.264 for the reference model and 0.220 for
ours, and the AUC that separates rib pixels from other lung pixels by the amount removed is the same (0.671 and 0.672). Ours is the more
conservative suppressor: it removes less rib edge and much less of everything else. Read with the audit, ours keeps lesions better per unit of
rib suppression (absorption a third of the reference model's at about 60\,\% of its rib-edge removal), but a share of the reference model's
higher absorption may simply reflect stronger suppression. The benchmark says nothing about lesions directly: it is a measure of what is
removed, not of what is kept.
\begin{table}[htbp]
\centering\small
\caption{Rib benchmark on VinDr-RibCXR (metrics only): gradient-energy reduction on rib-boundary pixels and on other lung pixels, selectivity
and AUC.}
\label{tab:rib_benchmark}
\fittable{tables/rib_benchmark.tex}
\end{table}

\section{Real nodules}
Real nodules answer RQ3 only in part (Table~\ref{tab:nodule_cnr}). On JSRT (154 nodules) the contrast-to-noise ratio of the nodule box rises
by 0.041 in the soft image of ours and by 0.018 in that of the reference model; the difference, $+0.023$ $[0.008,0.041]$, grows to $+0.032$
$[0.014,0.052]$ for the 123 nodules with at least 15\,\% bone overlap, where the reference model lowers contrast. On Node21 (1,474 boxes) the
reference model raises contrast more ($+0.103$ against $+0.078$; difference $-0.025$ $[-0.033,-0.018]$), also for the 1,109 bone-overlapping
boxes ($-0.017$ $[-0.026,-0.009]$). Against the no-twin control ours raises contrast in all four comparisons ($+0.007$ to $+0.027$, Holm
$p<0.001$). The twin regulariser thus helps real nodules relative to the same model without it, consistently; whether ours beats the reference
model depends on the data set, and on Node21 the reference model's stronger suppression raises box contrast more.
\begin{table}[htbp]
\centering\small
\caption{Contrast-to-noise change of real nodules on JSRT and Node21 (soft-tissue image minus original), all lesions and lesions with at least
15\,\% bone overlap; image-cluster bootstrap intervals.}
\label{tab:nodule_cnr}
\fittable{tables/nodule_cnr.tex}
\end{table}

\section{Downstream classification}
On NIH ChestX-ray14 (Table~\ref{tab:downstream_cls}) the three-seed ensemble of DenseNet-121 classifiers reaches a macro-AUROC of 0.818 on the
original images and 0.817 on the soft-tissue images of ours ($-0.001$ $[-0.004,0.003]$, patient-cluster bootstrap). Per class the paired
DeLong tests find four differences after Holm correction: ours is lower for effusion, infiltration and pleural thickening ($-0.004$, $-0.007$,
$-0.013$) and higher for nodule ($+0.011$). Adding the soft image as a second channel helps a little ($+0.006$ $[0.003,0.008]$), and so does the
reference model's soft image on its own ($+0.004$ $[0.002,0.007]$); CLAHE lowers the macro-AUROC ($-0.004$ $[-0.006,-0.002]$). The soft image
of ours is no substitute for the original, but the nodule class, the one the method was designed for, moves in its favour.
\begin{table}[htbp]
\centering\small
\caption{Downstream classification on NIH ChestX-ray14: macro-AUROC of the three-seed ensemble, per-seed mean $\pm$ SD and the difference to
the original image (patient-cluster bootstrap).}
\label{tab:downstream_cls}
\fittable{tables/downstream_cls.tex}
\end{table}

\section{Downstream detection}
For detection (Table~\ref{tab:downstream_det}, Figure~\ref{fig:downstream_forest}) the change in CPM with the soft image of ours alone is
$+0.001$ $[-0.023,0.023]$ on Node21, $+0.006$ $[-0.009,0.020]$ on TBX11K and $-0.009$ $[-0.016,-0.002]$ on VinBigData. The three-channel input
(original, soft and bone) gains $+0.028$ $[0.009,0.047]$ on Node21 and $+0.008$ $[0.002,0.014]$ on VinBigData; the reference model's soft image
gains $+0.029$ $[0.016,0.043]$ on TBX11K. The bone-overlap subgroup, recall at two false positives per image for lesions with and without at
least 15\,\% bone overlap compared as a difference in differences against the original input, shows no effect whose interval excludes zero
(Node21 $-0.006$ $[-0.071,0.061]$, VinBigData $+0.009$ $[-0.034,0.053]$); on TBX11K it cannot be estimated, because every test box overlaps bone
by at least 20\,\%. H4$'$b therefore finds no benefit of the soft image used alone; the gains come from giving the detector the soft and bone
images in addition to the original.
\begin{table}[htbp]
\centering\small
\caption{Downstream detection on Node21, TBX11K and VinBigData: CPM and mAP@50 (mean $\pm$ SD over three seeds), differences to the original
image and the difference in differences (DiD) of recall at two false positives per image (image bootstrap).}
\label{tab:downstream_det}
\fittable{tables/downstream_det.tex}
\end{table}
\begin{figure}[htbp]
\centering
\safefig{figures/downstream_forest.pdf}
\caption{Downstream effects of each input against the original image with 95\,\% intervals: macro-AUROC on NIH ChestX-ray14, CPM on the
detection sets.}
\label{fig:downstream_forest}
\end{figure}

\section{Comparison with the reference model and the controls}
The reference model is compared as released: a SwinV2-B suppressor that predicts at $1024$~px, trained by its authors on far more samples than
ours, run with its own pre- and post-processing. Our models predict at $512$~px with a tiny encoder. Differences between the two therefore mix
the effect of the regulariser with differences in resolution, capacity, data and training length; the matched no-twin control, which shares
everything with ours except the regulariser, is the comparison that isolates the regulariser, and it absorbs slightly more than the reference
model. The controls frame the scale of the effects: the identity removes nothing and absorbs nothing, CLAHE changes contrast everywhere and
does worse than the original image downstream, and homomorphic filtering removes low-frequency content without separating bone from soft
tissue. BS-LDM, trained on real dual-energy pairs from the test hospital, is far ahead in fidelity on that hospital's images and far behind
in the development audit; it answers a different question and is not part of a pre-specified comparison.
